\subsection{代数扩张中的赋值环}

命题 6. 设 K 为域，v 是 K 上的赋值，A 是其环，L 是 K 的代数扩张，A' 是 A 在 L 中的整闭包。设 V 为 L 上延拓 v 的赋值环的集合，M' 为 A' 的极大理想集合。则映射 V \(\mapsto\) m(V) \(\cap\) A' 是从 V 到 M' 的双射，而 \(m' \mapsto A'_{m'}\) 是逆双射。

A' 中的每个极大理想 \(m'\)（属于 \(A'\)）都满足 \(m' \cap A\) 是 A 的极大理想（第五章，第 2 节，第 1 节，命题 1），且 \(A_{m'}'\) 被 L 的某个赋值环 V 支配（因而它是 L 上延拓 v 的某个赋值的环）（§1，第 2 节，定理 2 的推论）。域 L 是由有限于 K 的 L 的子扩域 \(K_{\alpha}\) 组成的一个有向族的并；为了证明 \(V = A_{m'}'\)，只需证明 \(V \cap K_{\alpha} = A_{m'}' \cap K_{\alpha}\) 对所有 \(\alpha\) 都成立。现在，若记 \(A_{\alpha}' = A' \cap K_{\alpha}\)，则 \(A_{\alpha}'\) 是 A 在 \(K_{\alpha}\) 中的整闭包，因而是 \(K_{\alpha}\) 上延拓 v 的各赋值的环的交，而这些环 \(V_{i\alpha}\) 的数目有限，并且是局部环 \((\mathrm{A}_{\alpha}')_{\mathrm{m}'_{i\alpha}}\) 的 \(A_{\alpha}' (1 \leqslant i \leqslant n)\)，其中 \(m'_{i\alpha}\) 是 \(A_{\alpha}' (no. 3, Remark)\) 的不同极大理想；但 \(m' \cap A_{\alpha}'\) 是其中一个 \(m'_{i\alpha}\)，所以 \(V \cap K_{\alpha}\) 等于相应的局部环 \((\mathrm{A}'_{\alpha})_{m'_{i\alpha}} \subset \mathrm{A}'_{m'}\)，这就完成了 \(V = A'_{m'}\) 的证明。反之，若 \(V \in V\)，则 \(A' \subset V\)（§3，第 3 节，命题 6），且若 \(m' = m(V) \cap A'\)，则 \(m' \cap A = m\)，所以 \(m'\) 是 \(A'\) 的极大理想（第五章，第 2 节，第 1 节，命题 1），上述论证表明 \(V = A'_{m'}\)。

命题 7. 设 K 为域，L 是 K 的拟 Galois 扩张，f 和 \(f'\) 是取值于同一域 F 的 L 的位。假设 f 和 \(f'\) 在 K 上的限制相同。则存在 L 的一个 K-自同构 s，使得 \(f' = f \circ s\)。

设 A 为 K 上 f 与 \(f'\) 的公共限制所对应的位的环。f 与 \(f'\) 的环包含 A 在 L 中的整闭包 \(A'\)（§ 1，第 3 节，定理 3 的推论 3），因而（第五章，第 2 节，第 3 节，命题 6 的推论 1）存在 L 的一个 K-自同构 s，使得 \(f'\) 与 \(f \circ s\) 在 \(A'\) 上的限制相等；若 \(m'\) 是这些限制的公共核，则 \(m' \cap A\) 是 A 的极大理想，因而 \(m'\) 是 \(A'\) 的极大理想，且位 \(f'\) 与 \(f \circ s\) 在环 \(A_{m'}\) 上重合；但由命题 6，L 中支配 \(A_{m'}\) 的唯一赋值环就是环 \(A_{m'}\) 本身，所以位 \(f'\) 与 \(f \circ s\) 的环相同。

推论 1. 设 K 为域，v 是 K 上的赋值，L 是 K 的拟 Galois 扩张，v' 与 v'' 是 v 到 L 的两个延拓。则存在 L 的一个 K-自同构 s，使得 v'' 等价于 v' 与 s 的复合。

令 \(f'\) 和 \(f''\) 为分别与 \(v'\) 和 \(v''\) 相关的 K 的位；必要时以等价的位替换它们，可以设它们都将值取在 \(v\) 的剩余域的代数闭包中（第 1 节，命题 1）。于是存在 L 的一个 K-自同构 \(s\)，使得 \(f'' = f' \circ s\)（命题 7）；由位与赋值之间的对应（§ 3，第 3 节），\(v''\) 等价于 \(v' \circ s\)。

推论 2. 设 K 为域，f 是 K 的一个位（分别地，v 是 K 上的一个赋值），L 是 K 的纯不可分扩张。则 f（分别地，v）到 L 的所有延拓都等价。

L 是拟 Galois 扩张，且其唯一的自同构是恒等自同构。因此，推论 2 由命题 7（分别地，推论 1）得到。

命题 8. 设 K 为域，v 是 K 上的赋值，L 是 K 的次数为 n 的有限拟 Galois 扩张，且 \((v_{i}^{\prime})_{1\leq i\leq g}\) 是 v 到 L 的一个完备延拓系。则 \(e(v_{i}^{\prime}|v)\) 和 \(f(v_{i}^{\prime}|v)\) 的值 e、f 与 i 无关。于是 \(efg \leqslant n\)。若 v 的环 A 在 L 中的整闭包是有限生成的 A-模，则 \(efg = n\)。

这立即由定理 1（第 3 节）和定理 2（第 5 节）得到。

